%======================================================================
% Appendix A: Notation and the Greek alphabet
%======================================================================
\chapter{Notation and the Greek Alphabet}\label{app:notation}

This appendix collects the notation used in these notes, which follows the
conventions of most machine learning textbooks and research papers. The last
column says where each symbol is introduced.

\section*{Conventions}
\phantomsection
\addcontentsline{toc}{section}{Conventions}

\begin{itemize}
  \item Ordinary numbers (\emph{scalars}) are written in italic: $x$, $a$, $\lambda$.
  \item Vectors are bold lower-case letters, $\vx$, $\vw$, $\vtheta$, and are
    column vectors; their components are $x_1, \dots, x_n$.
  \item Matrices are bold capital letters, $\mA$, $\mW$, $\mSigma$, with entries
    $A_{ij}$ (row $i$, column $j$).
  \item Random variables are italic capitals, $X$, $Y$; their values are lower-case,
    $x$, $y$. Random vectors are bold capitals, $\vX$.
  \item A hat marks an estimate or a prediction: $\hat\theta$, $\hat y$.
  \item $\ln$ is the natural logarithm. In information theory, $\log$ means
    $\log_2$ when measuring in bits and $\ln$ when measuring in nats.
  \item Superscripts in parentheses number the layers of a network:
    $\mW^{(l)}$ is the weight matrix of layer $l$, not a power.
\end{itemize}

\section*{Symbols}
\phantomsection
\addcontentsline{toc}{section}{Symbols}

\renewcommand{\arraystretch}{1.25}
\begin{longtable}{@{}>{\raggedright\arraybackslash}p{3.6cm}
                   >{\raggedright\arraybackslash}p{7.3cm}
                   >{\raggedright\arraybackslash}p{3.0cm}@{}}
\toprule
\textbf{Symbol} & \textbf{Meaning} & \textbf{Introduced in}\\
\midrule
\endfirsthead
\toprule
\textbf{Symbol} & \textbf{Meaning} & \textbf{Introduced in}\\
\midrule
\endhead
\midrule
\multicolumn{3}{r}{\small\itshape continued on the next page}\\
\endfoot
\bottomrule
\endlastfoot
%
\multicolumn{3}{@{}l}{\textbf{\color{MLchapter}Numbers and sets}}\\
$\N$, $\Z$, $\Q$, $\R$ & natural numbers, integers, rational and real numbers & \cref{sec:kinds-of-numbers}\\
$x \in A$ & $x$ is an element of the set $A$ & \cref{sec:kinds-of-numbers}\\
$A \subset B$ & $A$ is a subset of $B$ & \cref{sec:kinds-of-numbers}\\
$\{x \in \R : x > 0\}$ & the set of all $x$ with the stated property & \cref{sec:kinds-of-numbers}\\
$[a, b]$, $(a, b)$ & closed and open intervals & \cref{sec:kinds-of-numbers}\\
$A \cup B$, $A \cap B$, $A^c$ & union, intersection, complement & \cref{sec:events}\\
$\emptyset$ & the empty set & \cref{sec:events}\\
$\abs{A}$ & number of elements of a finite set $A$ & \cref{sec:probability-rules}\\
$\R^n$ & the set of vectors with $n$ real components & \cref{sec:vectors}\\
\addlinespace
\multicolumn{3}{@{}l}{\textbf{\color{MLchapter}Arithmetic and algebra}}\\
$a^n$, $\sqrt a$, $a^{1/n}$ & powers and roots & \cref{sec:exponents}\\
$\sum_{i=1}^n x_i$ & the sum $x_1 + x_2 + \dots + x_n$ & \cref{sec:summation}\\
$\prod_{i=1}^n x_i$ & the product $x_1x_2\cdots x_n$ & \cref{sec:summation}\\
$n!$ & factorial, $1\times2\times\dots\times n$ & \cref{sec:summation}\\
$\binom nk$ & binomial coefficient, ``$n$ choose $k$'' & \cref{sec:probability-rules}\\
$\abs x$ & absolute value & \cref{sec:absolute-value}\\
$\approx$ & is approximately equal to & throughout\\
$\propto$ & is proportional to & \cref{sec:map}\\
\addlinespace
\multicolumn{3}{@{}l}{\textbf{\color{MLchapter}Functions}}\\
$f\colon A \to B$ & a function from $A$ to $B$ & \cref{sec:what-is-function}\\
$f \circ g$ & composition, $(f \circ g)(x) = f(g(x))$ & \cref{sec:composition}\\
$f^{-1}$ & inverse function & \cref{sec:inverse-functions}\\
$e$, $e^x = \exp(x)$ & Euler's number $2.718\ldots$, the exponential function & \cref{sec:exponential}\\
$\ln x$, $\log_bx$ & natural logarithm, logarithm to base $b$ & \cref{sec:logarithms}\\
$\sin$, $\cos$, $\tan$ & trigonometric functions & \cref{sec:trigonometry}\\
$\sigma(x)$ & sigmoid function $\frac{1}{1 + e^{-x}}$ (but see also $\sigma$ below) & \cref{sec:inverse-functions}\\
$\ReLU(x)$, $\tanh(x)$ & activation functions & \cref{sec:activation-gallery}\\
$\softmax(\vz)$ & softmax function & \cref{sec:softmax}\\
\addlinespace
\multicolumn{3}{@{}l}{\textbf{\color{MLchapter}Linear algebra}}\\
$\vx$, $x_i$ & a vector and its $i$-th component & \cref{sec:vectors}\\
$\mA$, $A_{ij}$ & a matrix and its entry in row $i$, column $j$ & \cref{sec:matrices}\\
$\vzero$, $\vone$ & the vector of zeros, the vector of ones & \cref{sec:vectors}\\
$\mI$ & identity matrix & \cref{sec:matrices}\\
$\vx\cdot\vy = \vx^\top\vy$ & dot product & \cref{sec:dot-product}\\
$\norm{\vx}$ & length (Euclidean or $L^2$ norm) & \cref{sec:vector-length}\\
$\norm{\vx}_1$, $\norm{\vx}_p$, $\norm{\vx}_\infty$ & $L^1$, $L^p$ and $L^\infty$ norms & \cref{sec:norms}\\
$\mA^\top$ & transpose & \cref{sec:transpose}\\
$\vx\vy^\top$ & outer product & \cref{sec:transpose}\\
$\mA^{-1}$ & inverse matrix & \cref{sec:inverses}\\
$\det\mA$ & determinant & \cref{sec:inverses}\\
$\spn(\vv_1, \dots, \vv_k)$ & span: all linear combinations & \cref{sec:span}\\
$\rank\mA$ & rank & \cref{sec:basis}\\
$\proj_{\va}(\vb)$ & projection of $\vb$ onto the line through $\va$ & \cref{sec:projections}\\
$\mA\vv = \lambda\vv$ & eigenvalue $\lambda$ and eigenvector $\vv$ & \cref{sec:eigen}\\
$\tr\mA$ & trace, the sum of the diagonal entries & \cref{sec:eigen}\\
$\mA = \mU\mSigma\mV^\top$ & singular value decomposition & \cref{sec:svd}\\
$\diag(d_1, \dots, d_n)$ & diagonal matrix with the given diagonal & \cref{sec:jacobian}\\
$\vu \odot \vv$ & componentwise (Hadamard) product & \cref{sec:matrix-calculus}\\
\addlinespace
\multicolumn{3}{@{}l}{\textbf{\color{MLchapter}Calculus}}\\
$\lim_{x \to a}f(x)$ & limit & \cref{sec:limits}\\
$f'(x)$, $\deriv fx$ & derivative & \cref{sec:derivative}\\
$f''(x)$ & second derivative & \cref{sec:extrema}\\
$\int_a^bf(x)\,\dd x$ & definite integral & \cref{sec:integration}\\
$\pderiv fx$ & partial derivative & \cref{sec:partial-derivatives}\\
$\nabla f$, $\nabla_{\vx}f$ & gradient (vector of partial derivatives) & \cref{sec:gradient}\\
$D_{\vu}f$ & directional derivative & \cref{sec:directional-derivatives}\\
$\mat J$ & Jacobian matrix & \cref{sec:jacobian}\\
$\mH$ & Hessian matrix & \cref{sec:hessian}\\
$\pderiv{L}{\mW}$ & matrix of the derivatives $\pderiv{L}{W_{ij}}$ & \cref{sec:matrix-calculus}\\
\addlinespace
\multicolumn{3}{@{}l}{\textbf{\color{MLchapter}Optimisation}}\\
$\vtheta$ & the parameters of a model & \cref{sec:optimisation-problems}\\
$L(\vtheta)$, $\ell$ & loss function, loss of one example & \cref{sec:optimisation-problems}\\
$\argmin_xf(x)$, $\argmax_xf(x)$ & the point where $f$ is smallest, largest & \cref{sec:optimisation-problems}\\
$\eta$ & learning rate & \cref{sec:gradient-descent}\\
$\lambda$ & regularisation strength; Lagrange multiplier & \cref{sec:norms,sec:lagrange}\\
\addlinespace
\multicolumn{3}{@{}l}{\textbf{\color{MLchapter}Probability}}\\
$\Omega$ & sample space & \cref{sec:events}\\
$P(A)$, $P(A \mid B)$ & probability, conditional probability & \cref{sec:probability-rules,sec:conditional}\\
$X$, $p(x)$, $F(x)$ & random variable, its PMF or PDF, its CDF & \cref{sec:random-variables}\\
$\E[X]$, $\mu$ & expected value (mean) & \cref{sec:expectation}\\
$\Var(X)$, $\sigma^2$, $\sigma$ & variance, standard deviation & \cref{sec:variance}\\
$\Cov(X, Y)$, $\Corr(X, Y)$, $\rho$ & covariance, correlation & \cref{sec:covariance}\\
$\mSigma$ & covariance matrix & \cref{sec:covariance}\\
$X \sim \mathcal{N}(\mu, \sigma^2)$ & $X$ has the normal distribution with mean $\mu$, variance $\sigma^2$ & \cref{sec:distributions}\\
$\operatorname{Bin}(n, p)$ & binomial distribution & \cref{sec:distributions}\\
i.i.d. & independent and identically distributed & \cref{sec:independence-events}\\
\addlinespace
\multicolumn{3}{@{}l}{\textbf{\color{MLchapter}Statistics}}\\
$\bar x$, $s^2$ & sample mean, sample variance & \cref{sec:descriptive}\\
$\hat\theta$ & estimate of the parameter $\theta$ & \cref{sec:estimators}\\
$\Bias$, $\MSE$, $\operatorname{SE}$ & bias, mean squared error, standard error & \cref{sec:estimators}\\
$\mathcal{L}(\theta)$ & likelihood & \cref{sec:mle}\\
$\mathcal{D}$ & a data set & \cref{sec:map}\\
$H_0$, $H_1$ & null and alternative hypotheses & \cref{sec:hypothesis-testing}\\
\addlinespace
\multicolumn{3}{@{}l}{\textbf{\color{MLchapter}Information theory and deep learning}}\\
$H(\vp)$ & entropy & \cref{sec:information-theory}\\
$H(\vp, \vq)$ & cross-entropy & \cref{sec:information-theory}\\
$D_{\KL}(\vp\,\Vert\,\vq)$ & Kullback--Leibler divergence & \cref{sec:information-theory}\\
$\mW^{(l)}$, $\vb^{(l)}$ & weights and biases of layer $l$ & \cref{sec:neural-networks}\\
$\vz^{(l)}$, $\va^{(l)}$ & pre-activations and activations of layer $l$ & \cref{sec:neural-networks}\\
$\varphi$ & activation function & \cref{sec:neural-networks}\\
$\vq$, $\vy$ & predicted probabilities, one-hot label & \cref{sec:softmax-gradient}\\
$\vdelta^{(l)}$ & error signal $\pderiv{L}{\vz^{(l)}}$ & \cref{sec:backpropagation}\\
$\mQ$, $\mK$, $\mV$ & queries, keys and values in attention & \cref{sec:attention}\\
\end{longtable}

\clearpage
\section*{The Greek alphabet}
\phantomsection
\addcontentsline{toc}{section}{The Greek alphabet}

Mathematics uses Greek letters as additional symbols. The table lists all of them,
with their typical meanings in machine learning. Letters that look exactly like
Latin ones (such as capital alpha, A) are not used as separate symbols.

\begin{center}
\renewcommand{\arraystretch}{1.2}
\begin{tabular}{@{}cclp{8.4cm}@{}}
\toprule
\textbf{Lower} & \textbf{Upper} & \textbf{Name} & \textbf{Typical uses}\\
\midrule
$\alpha$ & A & alpha & significance level; attention weights; step sizes\\
$\beta$ & B & beta & momentum coefficients $\beta_1$, $\beta_2$; regression coefficients\\
$\gamma$ & $\Gamma$ & gamma & discount factor in reinforcement learning; the gamma function\\
$\delta$ & $\Delta$ & delta & a small change ($\Delta x$); error signals $\vdelta$\\
$\epsilon$, $\varepsilon$ & E & epsilon & a small positive number; noise\\
$\zeta$ & Z & zeta & occasionally a variable\\
$\eta$ & H & eta & learning rate\\
$\theta$ & $\Theta$ & theta & parameters of a model; angles\\
$\iota$ & I & iota & rarely used\\
$\kappa$ & K & kappa & condition number\\
$\lambda$ & $\Lambda$ & lambda & eigenvalues ($\mLambda$: diagonal matrix of them); regularisation strength; rates\\
$\mu$ & M & mu & mean\\
$\nu$ & N & nu & degrees of freedom\\
$\xi$ & $\Xi$ & xi & random variables; slack variables\\
$o$ & O & omicron & not used (looks like the letter o)\\
$\pi$ & $\Pi$ & pi & the number $3.14159\ldots$; a policy; $\Pi$: product\\
$\rho$ & P & rho & correlation\\
$\sigma$ & $\Sigma$ & sigma & standard deviation; sigmoid; $\Sigma$: sum, covariance matrix\\
$\tau$ & T & tau & temperature; thresholds; time\\
$\upsilon$ & $\Upsilon$ & upsilon & rarely used\\
$\phi$, $\varphi$ & $\Phi$ & phi & activation functions; feature maps; $\Phi$: normal CDF\\
$\chi$ & X & chi & the $\chi^2$ (chi-squared) distribution\\
$\psi$ & $\Psi$ & psi & feature maps; parameters\\
$\omega$ & $\Omega$ & omega & frequencies; $\Omega$: sample space\\
\bottomrule
\end{tabular}
\end{center}
